% !TEX program = lualatex
\documentclass[aspectratio=169,professionalfonts]{beamer}
\newcommand{\slidemode}{1}  % カラー投影モード
\usepackage{beamer_template}
\usepackage{enumerate}

\newcommand{\LectureNum}{9}
\newcommand{\LectureTitle}{一般の周期関数のフーリエ級数（１）}
\newcommand{\CourseName}{応用数学}
\renewcommand{\TermName}{前期}

\begin{document}

\begin{frame}{第9回　一般の周期関数のフーリエ級数（１）}
  \begin{block}{本回の内容}
  一般の周期関数のフーリエ級数（１）
  \end{block}
  \vfill\hfill{\scriptsize \textcolor{gray}{第3章 フーリエ解析　§1}}
\end{frame}

\begin{frame}{定義：フーリエ展開（周期 T = 2l）}
  \begin{block}{}
  \[
    f(t) = a_{0}/2 + \sum_{n=1}^\infty (a_{n}\cos(n\pi t/l) + b_{n}\sin(n\pi t/l))
  \]
  \medskip\textbf{フーリエ係数}：
  \begin{align*}
    a_{0} &= (1/l) \int_{-l}^{l} f(t) dt \\
    a_{n} &= (1/l) \int_{-l}^{l} f(t)\cos(n\pi t/l) dt  (n = 1, 2, ...) \\
    b_{n} &= (1/l) \int_{-l}^{l} f(t)\sin(n\pi t/l) dt  (n = 1, 2, ...)
  \end{align*}
  \end{block}
\end{frame}

\begin{frame}{例題3　（p.\,83）}
  \begin{exampleblock}{問題}
    $f(x) = |x| (-1 \leq x < 1), f(x+2) = f(x)$ のフーリエ展開を求めよ\\[2pt]
  \end{exampleblock}
\end{frame}

\begin{frame}{例題3　解答}
  $f(x)$ は偶関数なので $b_{n} = 0,$ フーリエ余弦展開を求めればよい\\[4pt]
  $c_{0} = \int_{-1}^{1} |x| dx = 1$\\[4pt]
  $a_{n} = 2\int_{0}^{1} x \cos(n\pi x) dx = -2(1-(-1)^n)/(n^{2}\pi^{2})$
  \medskip\textbf{答}：$f(x) = 1/2 - (4/\pi^{2})(\cos \pi x + \cos(3\pi x)/9 + \cos(5\pi x)/25 + \ldots)$
  {\footnotesize \textcolor{gray}{\textit{$N=2, N=5$ での収束グラフあり（ p.83 ）}}}
\end{frame}

\begin{frame}{演習問題}
  \begin{exampleblock}{自分で解いてみよう}
  \begin{description}
    \item[問3] 次の関数のフーリエ級数を求めよ
    $(1) f(x) = { 1  (-1 \leq x < 0),  0  (0 \leq x < 1) },  f(x+2) = f(x)$\\
    $(2) g(x) = { 2  (-2 \leq x < 0),  4  (0 \leq x < 2) },  g(x+4) = g(x)$\\
    $(3) h(x) = { 2x+1  (-1/2 \leq x < 0),  0  (0 \leq x < 1/2) },  h(x+1) = h(x)$\\
  \end{description}
  \end{exampleblock}
\end{frame}

\end{document}
